\section{Trajectory Case Study: \texttt{aime\_1991\_p9}}
\label{sec:case-study}
\label{app:case-study-aime-1991-p9}

This appendix examines the complete, successful search trajectory for the theorem \texttt{aime\_1991\_p9} from the \texttt{miniF2F-valid} split. Proving this Olympiad-level trigonometry and algebra problem requires establishing:
\begin{align*}
\forall (x : \mathbb{R}) (m : \mathbb{Q}), \quad 
&\left( \frac{1}{\cos x} + \tan x = \frac{22}{7} \right) \land \left( \frac{1}{\sin x} + \cot x = m \right) \\
&\implies m.\mathrm{den} + m.\mathrm{num} = 44
\end{align*}
where $m \in \mathbb{Q}$ is the rational target, represented as the fraction $m = \frac{29}{15}$, so that $m.\mathrm{num} = 29$ and $m.\mathrm{den} = 15$. 

The theorem was successfully solved by GammaZero at a search depth of 7, generating 506 actions,
making 916 Lean verification calls, and containing 545 nodes (39 proof-state OR-nodes and 506 action
AND-nodes) in the final search graph. Below we trace the exact nested skeleton structure, the
coordination of local proofs, and present the complete reconstructed Lean 4 proof.

\subsection{Hierarchical Search Depth-7 Tree Topology}

To illustrate how GammaZero's hierarchical search operates, we map the exact path of the committed
proof tree. For context, the key mathematical goals associated with each state ID (\(s_i\)) in the tree
are:
\begin{itemize}
    \item \(s_0\): Root goal \((\uparrow m.\mathrm{den} + m.\mathrm{num} = 44)\).
    \item \(s_1\): Intermediate half-angle tangent landmark \((\tan(x/2) = 15/29)\).
    \item \(s_2\): Main trigonometric equation transformed to the half-angle parameter.
    \item \(s_3\), \(s_{13}\), \(s_{14}\): Rational parameterizations of \(\cos x\), \(\tan x\), and their algebraic combination.
    \item \(s_4\): Double-angle expansion for \(\cos x\).
    \item \(s_5\), \(s_6\), \(s_7\), \(s_9\): Derivation of the cosine-square substitution using secant and tangent-square identities \((1 + \tan^2(x/2) = 1/\cos^2(x/2))\).
    \item \(s_8\), \(s_{11}\): Critical non-degeneracy conditions \((\cos^2(x/2) \neq 0\) and \(\cos(x/2) \neq 0)\) to verify division validity.
    \item \(s_{21}\): Solving the rational algebraic equation \(\frac{1+t}{1-t} = \frac{22}{7} \implies t = \frac{15}{29}\).
    \item \(s_{22}\), \(s_{23}\), \(s_{27}\), \(s_{28}\), \(s_{29}\), \(s_{30}\), \(s_{31}\): Recursive scaffolding to establish the target cotangent/cosecant identity \((1/\sin x + \cot x = 1/\tan(x/2) = m = 29/15)\).
    \item \(s_{33}\), \(s_{34}\): Final extraction of the numerator \((m.\mathrm{num} = 29)\) and denominator \((m.\mathrm{den} = 15)\) of the rational result.
\end{itemize}

The tree below shows the compact logical hierarchy of how the goals are decomposed and solved:

\begin{lstlisting}[basicstyle=\scriptsize\ttfamily,columns=fullflexible,keepspaces=true,breaklines=true,showstringspaces=false]
D0 s0  final rational-denominator goal
`- a32 skeleton
   |- D1 s1  half-angle tangent value
   |  `- a65 skeleton
   |     |- D2 s2  transform first hypothesis to half-angle form
   |     |  `- a98 skeleton
   |     |     |- D3 s3  cosine half-angle rational form
   |     |     |  `- a133 skeleton
   |     |     |     |- D4 s4  cosine double-angle identity -> a134
   |     |     |     |- D4 s5  cosine-square denominator form
   |     |     |     |  `- a170 skeleton
   |     |     |     |     |- D5 s6  secant-square identity
   |     |     |     |     |  `- a211 skeleton
   |     |     |     |     |     |- D6 s7  tan-square identity
   |     |     |     |     |     |  `- a258 skeleton
   |     |     |     |     |     |     |- D7 s8  nonzero cosine-square -> a265
   |     |     |     |     |     |     `- D7 s9  tangent-square algebra -> a289
   |     |     |     |     |     `- D6 s10 unfold t definition -> a292
   |     |     |     |     `- D5 s11 nonzero half-cosine -> a301
   |     |     |     `- D4 s12 sine-square denominator form -> a324
   |     |     |- D3 s13 tangent double-angle identity -> a327
   |     |     `- D3 s14 combine reciprocal-cosine and tangent -> a339
   |     `- D2 s21 solve algebraic equation for t -> a344
   |- D1 s22 rational value of m
   |  `- a389 skeleton
   |     `- D2 s23 transform second hypothesis
   |        `- a423 skeleton
   |           |- D3 s27 cotangent-to-ratio identity -> a424
   |           `- D3 s28 half-angle reciprocal identity
   |              `- a476 skeleton
   |                 |- D4 s29 numerator half-angle identity -> a477
   |                 |- D4 s30 sine double-angle identity -> a485
   |                 `- D4 s31 reciprocal tangent identity -> a487
   |- D1 s33 numerator of m -> a495
   `- D1 s34 denominator of m -> a499
\end{lstlisting}

\subsection{Search Efficiency}

The search expanded 36 candidate actions at the root state (\texttt{state\_0}) and at the half-angle
state (\texttt{state\_1}) before finding the committed skeleton decomposition. This cost reflects two
sources of exploration.

\begin{itemize}
    \item \textbf{Syntactic and strategic variation:} the model tried alternative proof paths, including
    direct trigonometric identities, algebraic factorizations, and half-angle sine parameterizations.
    \item \textbf{Lean elaboration constraints:} many candidate skeletons failed quickly during verification
    because of type-checking errors, invalid bindings, or missing side conditions.
\end{itemize}

The priority scoring heuristic bounded this exploration: the run reached depth 7 while keeping the
active search focused on states that continued to expose useful scaffold structure.

\subsection{Local Tactic Coordination and Non-Degeneracy Proofs}

While skeleton actions establish the global logical hierarchy, local tactic actions prove the leaf subgoals:
\begin{itemize}
    \item \textbf{Algebraic Solvers:} State \texttt{state\_21} was closed by \texttt{action\_344} using a combination of \texttt{field\_simp} and \texttt{linarith} in 6 lines of code. \texttt{state\_14} was proved by \texttt{action\_339} using \texttt{by\_cases h : 1 + t = 0} to handle potential division-by-zero singularities.
    \item \textbf{Non-Degeneracy Proofs by Contradiction:} The critical division-by-zero prevention subgoals (\texttt{state\_8} and \texttt{state\_11}) were solved by \texttt{action\_265} and \texttt{action\_301} using contradiction. For instance, \texttt{action\_301} introduced the hypothesis $h_{\mathrm{cos\_half}} : \cos(x/2) = 0$, derived $\cos x = -1$ and $\tan x = 0$ via double-angle identities, and showed that this contradicts the main hypothesis $1/\cos x + \tan x = 22/7$ (since $-1 + 0 = -1 \neq 22/7$).
\end{itemize}

\subsection{Complete Reconstructed Lean 4 Proof}

Below is the complete, 114-line formal proof of \texttt{aime\_1991\_p9} reconstructed and verified by the Lean 4 compiler. Note that since the surrounding parent theorem opens the namespaces via \texttt{open BigOperators Real Nat Rat Finset Topology}, the identifier \texttt{cos} and \texttt{Real.cos} are fully interchangeable and resolve cleanly to the same mathematical function during Lean's elaboration. Similarly, the lemma \texttt{pow\_eq\_zero} successfully resolves under this context to show that $a^n = 0 \implies a = 0$ for natural $n > 0$.

\begin{lstlisting}[basicstyle=\small\fontspec{DejaVu Sans Mono},columns=fullflexible,keepspaces=true,breaklines=true,showstringspaces=false]
theorem aime_1991_p9
  (x : ℝ) (m : ℚ)
  (h₀ : 1/Real.cos x + Real.tan x = 22/7)
  (h₁ : 1/Real.sin x + Real.cot x = m) :
  ↑m.den + m.num = 44 := by
  have h_tan_half : Real.tan (x / 2) = 15 / 29 := by
    have h_trig_identity : 1 / Real.cos x + Real.tan x = (1 + Real.tan (x / 2)) / (1 - Real.tan (x / 2)) := by
      let t := Real.tan (x / 2)
      have h_cos : Real.cos x = (1 - t^2) / (1 + t^2) := by
        have h_cos_double : Real.cos x = (Real.cos (x / 2))^2 - (Real.sin (x / 2))^2 := by
          nth_rewrite 1 [← mul_div_cancel₀ x two_ne_zero]
          rw [two_mul, Real.cos_add]
          ring
        have h_cos_sq_val : (Real.cos (x / 2))^2 = 1 / (1 + t^2) := by
          have h_tan_id : 1 + t^2 = 1 / (Real.cos (x / 2))^2 := by
            have h_tan_sq_id : 1 + (Real.tan (x / 2))^2 = 1 / (Real.cos (x / 2))^2 := by
              have h_sec_sq_id : (Real.cos (x / 2))^2 ≠ 0 := by
                intro h_cos_sq_zero
                have h_cos_half_zero : cos (x / 2) = 0 := pow_eq_zero h_cos_sq_zero
                have h_x_double : x = 2 * (x / 2) := by ring
                have h_cos_x : cos x = -1 := by
                  rw [h_x_double, cos_two_mul, h_cos_half_zero]
                  simp [sin_sq]
                have h_tan_x : tan x = 0 := by
                  rw [tan_eq_sin_div_cos, h_x_double, sin_two_mul, h_cos_half_zero]
                  simp
                rw [h_cos_x, h_tan_x] at h₀
                norm_num at h₀
              have h_standard_identity : (Real.tan (x / 2))^2 + 1 = 1 / (Real.cos (x / 2))^2 := by
                have h_cos_nz : Real.cos (x / 2) ≠ 0 := by
                  intro h_zero
                  apply h_sec_sq_id
                  rw [h_zero, zero_pow (by norm_num)]
                rw [Real.tan_eq_sin_div_cos, div_pow]
                field_simp [h_cos_nz]
                rw [Real.sin_sq_add_cos_sq]
              rw [add_comm]
              exact h_standard_identity
            have h_t_def : t = Real.tan (x / 2) := by
              rfl
            rw [h_t_def, h_tan_sq_id]
          have h_cos_nz : Real.cos (x / 2) ≠ 0 := by
            intro h_cos_half
            have h_cos_x : Real.cos x = -1 := by
              rw [show x = 2 * (x / 2) by ring, Real.cos_two_mul, h_cos_half]
              ring
            have h_sin_x : Real.sin x = 0 := by
              rw [show x = 2 * (x / 2) by ring, Real.sin_two_mul, h_cos_half]
              ring
            have h_tan_x : Real.tan x = 0 := by
              rw [Real.tan_eq_sin_div_cos, h_sin_x, h_cos_x]
              exact zero_div (-1)
            rw [h_cos_x, h_tan_x] at h₀
            norm_num at h₀
          field_simp [h_cos_nz] at h_tan_id ⊢
          exact h_tan_id
        have h_sin_sq_val : (Real.sin (x / 2))^2 = t^2 / (1 + t^2) := by
          have h_iden : (Real.sin (x / 2))^2 + (Real.cos (x / 2))^2 = 1 := sin_sq_add_cos_sq (x / 2)
          rw [h_cos_sq_val] at h_iden
          have h_nz : 1 + t^2 ≠ 0 := by nlinarith
          field_simp [h_nz] at h_iden ⊢
          linear_combination h_iden
        rw [h_cos_double, h_cos_sq_val, h_sin_sq_val]
        field_simp
      have h_tan : Real.tan x = (2 * t) / (1 - t^2) := by
        rw [show x = 2 * (x / 2) by ring, Real.tan_two_mul]
      have h_algebraic_step : 1 / ((1 - t^2) / (1 + t^2)) + (2 * t) / (1 - t^2) = (1 + t) / (1 - t) := by
        rw [one_div, inv_div, ← add_div]
        have h_sq : 1 + t^2 + 2 * t = (1 + t) * (1 + t) := by ring
        have h_diff : 1 - t^2 = (1 - t) * (1 + t) := by ring
        rw [h_sq, h_diff]
        by_cases h : 1 + t = 0
        · simp [h]
        · exact mul_div_mul_right (1 + t) (1 - t) h
      rw [h_cos, h_tan]
      exact h_algebraic_step
    have h_eq : (1 + Real.tan (x / 2)) / (1 - Real.tan (x / 2)) = 22 / 7 := by
      rw [← h_trig_identity]
      exact h₀
    have h_algebra : ∀ t : ℝ, (1 + t) / (1 - t) = 22 / 7 → t = 15 / 29 := by
      intro t ht
      have h_ne : 1 - t ≠ 0 := by
        intro h
        have : 1 + t = 2 := by linarith
        rw [h] at ht
        simp [div_zero] at ht
        norm_num at ht
      field_simp [h_ne] at ht
      linarith
    apply h_algebra (Real.tan (x / 2)) h_eq
  have h_m_val : m = 29 / 15 := by
    have h_trig_id : 1 / Real.sin x + Real.cot x = 1 / Real.tan (x / 2) := by
      have h_sum_frac : 1 / Real.sin x + Real.cot x = (1 + Real.cos x) / Real.sin x := by
        rw [Real.cot_eq_cos_div_sin, add_div]
      have h_half_angle : (1 + Real.cos x) / Real.sin x = 1 / Real.tan (x / 2) := by
        have h_num_exp : 1 + Real.cos x = 2 * Real.cos (x / 2) ^ 2 := by
          nth_rewrite 1 [show x = 2 * (x / 2) by ring]
          rw [Real.cos_two_mul, sq]
          ring
        have h_den_exp : Real.sin x = 2 * Real.sin (x / 2) * Real.cos (x / 2) := by
          rw [← Real.sin_two_mul, mul_div_cancel₀ x two_ne_zero]
        have h_tan_inv : 1 / Real.tan (x / 2) = Real.cos (x / 2) / Real.sin (x / 2) := by
          rw [tan_eq_sin_div_cos, one_div_div]
        rw [h_num_exp, h_den_exp, h_tan_inv]
        field_simp
      rw [h_sum_frac, h_half_angle]
    have h_m_real : (m : ℝ) = ((29 / 15 : ℚ) : ℝ) := by
      rw [← h₁, h_trig_id, h_tan_half]
      norm_num
    exact Rat.cast_inj.mp h_m_real
  have h_num : m.num = 29 := by
    rw [h_m_val]
    norm_num
  have h_den : m.den = 15 := by
    rw [h_m_val]
    norm_num
  rw [h_num, h_den]
  norm_num
\end{lstlisting}
